% GENERATED FILE. Edit canonical lessons, then run npm run generate:books.
% canonical-source-hash: fnv1a64:29f0eb8c18541f8b
% canonical-lessons: BN-C300-pathok, BN-C300-shrota, BN-C300-bigyopti, BN-C300-prochar-kora, BN-C300-songkhya

\chapter{Readers and Listeners}
\label{ch:bn-readers-and-listeners}
\hlchaptermodality{300}

\begin{chapteropening}
\textbf{By the end of this chapter:} \emph{I can talk about readers, listeners and public notices.}

Complete the last lesson of chapter 300: I can talk about readers, listeners and public notices.
\end{chapteropening}

\section[\texorpdfstring{\bn{পাঠক} (pāṭhôk)}{পাঠক (pāṭhôk)}]{\bn{পাঠক} (pāṭhôk) — a reader}

\label{lesson:BN-C300-pathok}



\begin{quote}
Before the new one: say the Bengali for live, then the Bengali for a viewer.
\end{quote}

\subsection*{You'll want to know: \bn{পাঠক}}
\textbf{\bn{পাঠক}} — \emph{pāṭhôk} — \textquotedblleft{}a reader\textquotedblright{}.

\begin{quote}
\bn{এই} \bn{পত্রিকার} \bn{অনেক} \bn{পাঠক} \bn{গ্রামে} \bn{থাকেন।}
\end{quote}

\begin{quote}
\emph{ei pôtrikār ônek pāṭhôk grāme thāken.}
\end{quote}

\begin{quote}
Many of this magazine's readers live in villages.
\end{quote}

\begin{cousinweb}
\textbf{\bn{পাঠক}} (\emph{pāṭhôk}) is Sanskrit \emph{pāṭhaka}, 'a reciter, a reader' (Monier-Williams), from the same root as \textbf{\bn{পড়া}}, to read.
\end{cousinweb}

\begin{grammarlens}[title={What you've built}]
One more word for telling what happened.
\end{grammarlens}

\subsection*{Your turn}
\begin{itemize}
\raggedright
  \item \emph{Say it:} \emph{pāṭhôk}
  \item \emph{Say it:} \emph{pāṭhôk}, once more
  \item \emph{Say it:} \emph{ei pôtrikār ônek pāṭhôk grāme thāken.} — Many of this magazine's readers live in villages.
  \item \emph{Recall:} say the Bengali for live, then the Bengali for a viewer, then say \emph{pāṭhôk} again
\end{itemize}

\begin{tcolorbox}[breakable,colback=teal!4,colframe=teal!35!black,title={Before you move on}]
What does \bn{পাঠক} mean? (\textquotedblleft{}A reader\textquotedblright{}.) Say it once more.
\end{tcolorbox}
\section[\texorpdfstring{\bn{শ্রোতা} (shrotā)}{শ্রোতা (shrotā)}]{\bn{শ্রোতা} (shrotā) — a listener}

\label{lesson:BN-C300-shrota}



\begin{quote}
Before the new one: say the Bengali for a viewer, then the Bengali for a reader.
\end{quote}

\subsection*{You'll want to know: \bn{শ্রোতা}}
\textbf{\bn{শ্রোতা}} — \emph{shrotā} — \textquotedblleft{}a listener\textquotedblright{}.

\begin{quote}
\bn{রেডিওর} \bn{শ্রোতারা} \bn{ফোন} \bn{করে} \bn{প্রশ্ন} \bn{জিজ্ঞাসা} \bn{করলেন।}
\end{quote}

\begin{quote}
\emph{reḍior shrotārā phon kôre prôshno jigyāsā kôrlen.}
\end{quote}

\begin{quote}
The radio listeners phoned in and asked questions.
\end{quote}

\begin{cousinweb}
\textbf{\bn{শ্রোতা}} (\emph{shrotā}) is Sanskrit \emph{śrotṛ}, 'one who hears' (Monier-Williams), in its nominative \emph{śrotā}; it belongs with \textbf{\bn{শোনা}}, to hear.
\end{cousinweb}

\begin{grammarlens}[title={What you've built}]
One more word for telling what happened.
\end{grammarlens}

\subsection*{Your turn}
\begin{itemize}
\raggedright
  \item \emph{Say it:} \emph{shrotā}
  \item \emph{Say it:} \emph{shrotā}, once more
  \item \emph{Say it:} \emph{reḍior shrotārā phon kôre prôshno jigyāsā kôrlen.} — The radio listeners phoned in and asked questions.
  \item \emph{Recall:} say the Bengali for a viewer, then the Bengali for a reader, then say \emph{shrotā} again
\end{itemize}

\begin{tcolorbox}[breakable,colback=teal!4,colframe=teal!35!black,title={Before you move on}]
What does \bn{শ্রোতা} mean? (\textquotedblleft{}A listener\textquotedblright{}.) Say it once more.
\end{tcolorbox}
\section[\texorpdfstring{\bn{বিজ্ঞপ্তি} (bigyôpti)}{বিজ্ঞপ্তি (bigyôpti)}]{\bn{বিজ্ঞপ্তি} (bigyôpti) — an official notice; an announcement}

\label{lesson:BN-C300-bigyopti}



\begin{quote}
Before the new one: say the Bengali for a reader, then the Bengali for a listener.
\end{quote}

\subsection*{You'll want to know: \bn{বিজ্ঞপ্তি}}
\textbf{\bn{বিজ্ঞপ্তি}} — \emph{bigyôpti} — \textquotedblleft{}an official notice; an announcement\textquotedblright{}.

\begin{quote}
\bn{স্টেশনে} \bn{একটা} \bn{বিজ্ঞপ্তি} \bn{ছিল} \bn{যে} \bn{ট্রেন} \bn{দেরিতে} \bn{আসবে।}
\end{quote}

\begin{quote}
\emph{sṭeshône ekṭā bigyôpti chhilô je ṭren derite āsbe.}
\end{quote}

\begin{quote}
There was a notice at the station saying the train would be late.
\end{quote}

\begin{cousinweb}
\textbf{\bn{বিজ্ঞপ্তি}} (\emph{bigyôpti}) is Sanskrit \emph{vijñapti}, 'information, a report' (Monier-Williams). Its \bn{জ্ঞ} is read \emph{gy}, as in \textbf{\bn{জ্ঞান}}.
\end{cousinweb}

\begin{grammarlens}[title={What you've built}]
One more word for telling what happened.
\end{grammarlens}

\subsection*{Your turn}
\begin{itemize}
\raggedright
  \item \emph{Say it:} \emph{bigyôpti}
  \item \emph{Say it:} \emph{bigyôpti}, once more
  \item \emph{Say it:} \emph{sṭeshône ekṭā bigyôpti chhilô je ṭren derite āsbe.} — There was a notice at the station saying the train would be late.
  \item \emph{Recall:} say the Bengali for a reader, then the Bengali for a listener, then say \emph{bigyôpti} again
\end{itemize}

\begin{tcolorbox}[breakable,colback=teal!4,colframe=teal!35!black,title={Before you move on}]
What does \bn{বিজ্ঞপ্তি} mean? (\textquotedblleft{}An official notice; an announcement\textquotedblright{}.) Say it once more.
\end{tcolorbox}
\section[\texorpdfstring{\bn{প্রচার} \bn{করা} (prôchār kôrā)}{প্রচার করা (prôchār kôrā)}]{\bn{প্রচার} \bn{করা} (prôchār kôrā) — to publicize; to spread the word}

\label{lesson:BN-C300-prochar-kora}



\begin{quote}
Before the new one: say the Bengali for a listener, then the Bengali for an official notice.
\end{quote}

\subsection*{You'll want to know: \bn{প্রচার} \bn{করা}}
\textbf{\bn{প্রচার} \bn{করা}} — \emph{prôchār kôrā} — \textquotedblleft{}to publicize; to spread the word\textquotedblright{}.

\begin{quote}
\bn{স্বেচ্ছাসেবকরা} \bn{গ্রামে} \bn{গ্রামে} \bn{খবরটা} \bn{প্রচার} \bn{করল।}
\end{quote}

\begin{quote}
\emph{svechchhāsebôkrā grāme grāme khôbôrṭā prôchār kôrlô.}
\end{quote}

\begin{quote}
The volunteers spread the word from village to village.
\end{quote}

Past (he/she …): \bn{তিনি} \bn{বিজ্ঞাপনটা} \bn{প্রচার} \bn{করলেন}, \emph{tini bigyāpônṭā prôchār kôrlen}.

\begin{cousinweb}
\textbf{\bn{প্রচার} \bn{করা}} (\emph{prôchār kôrā}) puts Sanskrit \emph{pracāra}, 'a going forth, a coming into view' (Monier-Williams), before \bn{করা}: to send word out among people.
\end{cousinweb}

\begin{grammarlens}[title={What you've built}]
One more word for telling what happened.
\end{grammarlens}

\subsection*{Your turn}
\begin{itemize}
\raggedright
  \item \emph{Say it:} \emph{prôchār kôrā}
  \item \emph{Say it:} \emph{prôchār kôrā}, once more
  \item \emph{Say it:} \emph{svechchhāsebôkrā grāme grāme khôbôrṭā prôchār kôrlô.} — The volunteers spread the word from village to village.
  \item \emph{Recall:} say the Bengali for a listener, then the Bengali for an official notice, then say \emph{prôchār kôrā} again
\end{itemize}

\begin{tcolorbox}[breakable,colback=teal!4,colframe=teal!35!black,title={Before you move on}]
What does \bn{প্রচার} \bn{করা} mean? (\textquotedblleft{}To publicize; to spread the word\textquotedblright{}.) Say it once more.
\end{tcolorbox}
\section[\texorpdfstring{\bn{সংখ্যা} (sôngkhyā)}{সংখ্যা (sôngkhyā)}]{\bn{সংখ্যা} (sôngkhyā) — an issue (of a paper); a number}

\label{lesson:BN-C300-songkhya}



\begin{quote}
Before the new one: say the Bengali for an official notice, then the Bengali for to publicize.
\end{quote}

\subsection*{You'll want to know: \bn{সংখ্যা}}
\textbf{\bn{সংখ্যা}} — \emph{sôngkhyā} — \textquotedblleft{}an issue (of a paper); a number\textquotedblright{}.

\begin{quote}
\bn{পত্রিকার} \bn{নতুন} \bn{সংখ্যায়} \bn{একটা} \bn{লম্বা} \bn{গল্প} \bn{আছে।}
\end{quote}

\begin{quote}
\emph{pôtrikār nôtun sôngkhyāy ekṭā lômbā gôlpo āchhe.}
\end{quote}

\begin{quote}
The new issue of the magazine has a long story in it.
\end{quote}

\begin{cousinweb}
\textbf{\bn{সংখ্যা}} (\emph{sôngkhyā}) is Sanskrit \emph{saṃkhyā}, 'a number, a reckoning' (Monier-Williams); a magazine's issues are counted, so each is a \bn{সংখ্যা}.
\end{cousinweb}

\begin{grammarlens}[title={What you've built}]
One more word for telling what happened.
\end{grammarlens}

\subsection*{Your turn}
\begin{itemize}
\raggedright
  \item \emph{Say it:} \emph{sôngkhyā}
  \item \emph{Say it:} \emph{sôngkhyā}, once more
  \item \emph{Say it:} \emph{pôtrikār nôtun sôngkhyāy ekṭā lômbā gôlpo āchhe.} — The new issue of the magazine has a long story in it.
  \item \emph{Recall:} say the Bengali for an official notice, then the Bengali for to publicize, then say \emph{sôngkhyā} again
\end{itemize}

\begin{tcolorbox}[breakable,colback=teal!4,colframe=teal!35!black,title={Before you move on}]
What does \bn{সংখ্যা} mean? (\textquotedblleft{}An issue (of a paper); a number\textquotedblright{}.) Say it once more.
\end{tcolorbox}
